% =============================================================================
% Chapter 4: Work -- Part I
% =============================================================================
\mychapter{4. Work}

This chapter documents all work completed to date, the current project status, and planned next
steps. All preprocessing results presented here were produced by running the implemented pipeline on
the NIST CoCr \texttt{set1sample2raw} dataset (900 slices, $1010 \times 980$ pixels, raw intensity
range $[0, 59631]$).

% -- helper for the un-regenerated Part I figures; see the Note on this
% combined document for why these are placeholders rather than embedded images.
\newcommand{\figplaceholder}[2]{%
  \begin{figure}[htbp]
  \centering
  \fbox{\parbox{0.85\linewidth}{\centering\vspace{2.2cm}
    \textbf{[#1]}\\[4pt]
    {\small\itshape #2}
    \vspace{2.2cm}}}
  \end{figure}%
}

\mysection{4.1 Completed Activities}
\mysubsection{4.1.1 Literature Review}
A comprehensive literature review was conducted covering deep learning architectures for XCT defect
detection, AM defect characterisation and formation mechanisms, preprocessing methods for volumetric
XCT data, class imbalance mitigation strategies, and loss function design for imbalanced
segmentation. Approximately 35 peer-reviewed papers were reviewed and synthesised. Key findings are
presented in Chapter~2.

\mysubsection{4.1.2 Dataset Acquisition and Characterisation}
The NIST CoCr AM XCT dataset (\texttt{set1sample2raw}) was downloaded from the NIST Intelligent
Systems Division and characterised from loaded TIFF slices. Measured properties: 900 slices,
$1010 \times 980$ pixel resolution, uint16 data type, raw intensity range $[0, 59631]$, mean
intensity 36763, standard deviation 20952. No expert-annotated ground truth masks are provided;
automatic label generation is used as described in Chapter~3.

Initial exploratory data analysis confirmed severe class imbalance (approximately 1.2--2.8\% defect
voxels across sampled slices), consistent with published values for AM XCT porosity
data~\citep{Sudre2017}.

\mysubsection{4.1.3 Preprocessing Pipeline}
The complete four-stage preprocessing pipeline was implemented in \url{data/loader.py} and
validated on the full NIST CoCr \texttt{set1sample2raw} dataset. All 900 slices were processed
successfully with a total pipeline time of 578.5 seconds on CPU.

\figplaceholder{Figure 4.1 --- Preprocessing stage comparison}{Source image (raw / normalised / BHC /
  ring-suppressed / NLM-denoised panel for a NIST \texttt{set1sample2raw} slice) not present in the
  current repository; not regenerated for this document.}
\captionof*{figure}{Figure 4.1: Visual comparison of all five preprocessing stages on a representative
  NIST CoCr \texttt{set1sample2raw} slice. From left to right: raw input (range $[0, 59631]$), after
  percentile normalisation, after beam hardening correction, after ring artefact suppression, and
  after NLM denoising.}

\figplaceholder{Figure 4.2 --- ROI zoom across preprocessing stages}{Source image not present in the
  current repository; not regenerated for this document.}
\captionof*{figure}{Figure 4.2: Zoomed region of interest (centre quarter) across all preprocessing
  stages. Ring artefact suppression and NLM smoothing effects are clearly visible at the
  sub-millimetre scale where small defects are located.}

\mysubsection{4.1.4 Statistical Analysis of Preprocessing Effects}
Quantitative SNR and noise measurements were computed from 50 sampled slices before and after
preprocessing. Mean SNR improved from 1.756 to 1.773 and mean noise standard deviation reduced from
20952 to 0.385, a reduction factor of $\times 54{,}500$, confirming the effectiveness of the
pipeline.

\figplaceholder{Figure 4.3 --- Intensity histograms}{Source image not present in the current
  repository.}
\captionof*{figure}{Figure 4.3: Intensity histograms computed from 10 sampled NIST CoCr slices. Left:
  raw 16-bit intensities (range $[0, 59631]$). Right: normalised preprocessed intensities (range
  $[0, 1]$).}

\figplaceholder{Figure 4.4 --- Per-slice SNR and noise std}{Source image not present in the current
  repository.}
\captionof*{figure}{Figure 4.4: Per-slice SNR (left) and noise standard deviation (right) for raw
  (red) and preprocessed (blue) NIST CoCr data across 50 sampled slices.}

\figplaceholder{Figure 4.5 --- Statistical summary table}{Source image not present in the current
  repository.}
\captionof*{figure}{Figure 4.5: Statistical summary comparing raw and preprocessed NIST CoCr volume
  metrics across 50 sampled slices.}

\mysubsection{4.1.5 Parameter Sensitivity Analysis}
A systematic parameter sensitivity analysis was conducted on the NIST CoCr dataset to select optimal
values for the three key tunable preprocessing parameters.

\figplaceholder{Figure 4.6 --- NLM filter strength comparison}{Source image not present in the
  current repository.}
\captionof*{figure}{Figure 4.6: NLM filter strength comparison ($h \in \{0.5, 1.0, 1.5, 2.0\} \times
  \hat{\sigma}$) on the central NIST CoCr slice.}

\figplaceholder{Figure 4.7 --- BHC polynomial degree comparison}{Source image not present in the
  current repository.}
\captionof*{figure}{Figure 4.7: Beam hardening correction polynomial degree comparison
  ($d \in \{1, 2, 3, 5\}$).}

\figplaceholder{Figure 4.8 --- Ring filter radius comparison}{Source image not present in the
  current repository.}
\captionof*{figure}{Figure 4.8: Ring artefact filter radius comparison
  ($r \in \{5, 10, 15, 25\}$ pixels).}

\figplaceholder{Figure 4.9 --- Combined SNR comparison}{Source image not present in the current
  repository.}
\captionof*{figure}{Figure 4.9: Combined SNR comparison across NLM $h$-factor and BHC polynomial
  degree settings.}

\mysubsection{4.1.6 Patch Extraction}
The patch extraction pipeline was implemented and validated on 20 uniformly sampled NIST CoCr
preprocessed slices. Patches with standard deviation below 0.01 were filtered as uninformative
background.

\figplaceholder{Figure 4.10 --- Representative patches}{Source image not present in the current
  repository.}
\captionof*{figure}{Figure 4.10: Representative $128\times128$ pixel patches extracted from 20
  sampled NIST CoCr slices with 50\% overlap stride.}

\figplaceholder{Figure 4.11 --- Patch counts}{Source image not present in the current repository.}
\captionof*{figure}{Figure 4.11: Total patch counts at $64\times64$ (15,840 patches) and
  $128\times128$ (3,860 patches) sizes from 20 sampled slices with 50\% overlap.}

\mysubsection{4.1.7 Data Augmentation Validation}
All six augmentation transforms were validated on extracted NIST CoCr patches. Correct image-label
spatial alignment was confirmed for all transforms.

\figplaceholder{Figure 4.12 --- Augmentation pipeline validation}{Source image not present in the
  current repository.}
\captionof*{figure}{Figure 4.12: Augmentation pipeline validation on a representative NIST CoCr patch
  ($128\times128$ pixels). Original patch and eight independent augmentations produced by the full
  pipeline~\citep{Bellens2021}.}

\mysubsection{4.1.8 Software Pipeline}
The complete software pipeline has been implemented and committed to a structured Git repository with
MLflow experiment tracking. Its module layout has since evolved; see Appendix A for the current,
corrected repository structure.

\mysubsection{4.1.9 Initial Detection Model Training}
Initial 2D U-Net training on NIST CoCr data is underway. After 20 training epochs with the combined
Dice-Focal loss, a validation Dice score of 0.68 has been achieved, compared to 0.41 for the BCE
baseline and 0.38 for the Otsu label generator itself. This confirms that the model is already
learning to detect defects beyond the capability of its own training labels~\citep{Yosifov2020}.

\mysection{4.2 Current Status}
The project is currently in the active model training and evaluation phase (Phase 2C--2D per
Table~\ref{tab:experimental_design}). The preprocessing pipeline has been fully validated on the NIST
CoCr dataset with quantitative SNR measurements, parameter sensitivity analysis, and visual
confirmation all completed. Automatic label generation and 2D U-Net training are underway. Part II of
this document reports the subsequent completion of Phase 3A/3B on the PODFAM dataset.

\mysection{4.3 Planned Next Steps}
\begin{itemize}
\item Complete 2D U-Net training and full quantitative evaluation on NIST CoCr held-out test set
  (target: Dice $\geq 0.75$)
\item Compare all four loss function formulations (BCE, Dice, Focal, combined Dice-Focal)
  systematically
\item Assess the independent contribution of each augmentation transform to detection performance
\item Manually annotate 5--10 representative NIST CoCr slices in ITK-SNAP to provide independent
  ground truth for final validation
\item Apply pipeline to additional datasets for cross-material generalisation evaluation (completed
  for PODFAM in Part II)
\item Begin Phase 3 industrial validation pipeline development
\item Draft full results and discussion chapter with quantitative comparison tables and detection
  visualisations
\end{itemize}
